\SetPractica{Calibración del material volumétrico}{Verificación de la calibración del material volumétrico}

\section{Introducción teórica}

El material volumétrico tiene por finalidad la medición exacta de volúmenes y debe ser controlado antes de utilizarlo. Para ello se requiere pesar la cantidad de agua pura (contenida en los matraces volumétricos) o transferida (por pipetas y buretas), a una temperatura dada, y calcular el volumen obtenido a partir de la masa pesada. Es importante que, antes de utilizar cualquier material volumétrico, se examine si las paredes del recipiente de medida están engrasadas. 

Para verificar esto se debe enjuagar el material con agua; cuando la superficie de vidrio está limpia y libre de grasa, el agua se extiende y deja una película invisible cuando se deja correr sobre ella. Si el agua no las humedece uniformemente, se debe limpiar.

Para la limpieza, muchas veces es suficiente una disolución de un detergente común. En caso de que no fuera suficiente, se puede utilizar mezcla crómica o una disolución de hidróxido de potasio en alcohol (esta última no debe dejarse mucho tiempo en contacto con el vidrio porque lo ataca lentamente). 

Siempre que se utilice una disolución de limpieza, el recipiente se lavará cuidadosamente, primero con agua corriente y después con agua destilada para verificar que las paredes queden uniformemente humedecidas. 

El material aforado no debe ser secado en estufa ya que puede provocar distorsión del vidrio y causar un cambio en el volumen

Para obtener la verificación del volumen calibrado a partir de la masa de agua es importante tener en cuenta que: 

\begin{enumerate}
    \item La densidad del agua varía con la temperatura 
    \item El volumen del recipiente de vidrio varía con la temperatura 
    \item El agua que llena el recipiente se pesa en aire 
\end{enumerate}

Cuando se verifica la calibración del material de vidrio se deben tomar en consideración estos factores para calcular el volumen contenido o vertido por el material a 20°C. 

Si trabajamos a temperatura ambiente (cercana a los 20°C) el segundo factor (el volumen del recipiente de vidrio varía con la temperatura) introduce correcciones muy pequeñas por lo que a efectos prácticos no lo consideraremos en el cálculo

\section{Objetivos}

\begin{itemize}
    \item 
\end{itemize}

\section{Materiales, equipos y reactivos}

\begin{table}[H]
\centering{\small\setstretch{1.25}
\begin{tabular}{|m{0.3\linewidth}|m{0.3\linewidth}|m{0.3\linewidth}|} \hline
    
    \thead{Equipos} & \thead{Materiales} & \thead{Reactivos} \\ \hline
    
    {Balanza analítica \par}
    {Termómetro}
    
    &
    
    {Vasos de precipitado \par}
    {Pipetas \par}
    {Perilla de succión}
    
    & 
    
    {Agua destilada \par}
    
    \\ \hline
\end{tabular}}\end{table}

\section{Procedimientos}

\subsection{Verificación e la calibración de pipeta volumétrica}
\begin{enumerate}
    \item Colocar agua destilada en un vaso de pp
    \item Medir la temperatura del agua.
    \item Llenar la pipeta de 5ml con agua destilada del vaso de pp
    \item Colocar el matraz volumétrico de 20ml limpio vacío con tapadera en la balanza analítica
    \item Tarar. 
    \item Quitar el tapón del matraz
    \item Transferir el agua al matraz aforado de 20ml
    \item Tapar el matraz
    \item Convertir la masa a volumen corregido (mL) de acuerdo con la temperatura
    \item Repetir este procedimiento 4 veces más. 
    \item Calcular las diferencias entre los volumen medido y los corregidos
    \item Calcular desviación estándar de los valores obtenidos
    \item Cotejar con la incertidumbre de la pipeta
    \item Determinar si es apta para usarse o no
\end{enumerate}

\subsection{Verificación de la bureta}
\begin{enumerate}
    \item Colocar agua destilada en un vaso de pp
    \item Medir la temperatura del agua.
    \item Llenar la bureta de 25ml con esta agua hasta la marca de 0
    \item Colocar el matraz volumétrico de 25ml limpio vacío con tapadera en la balanza analítica
    \item Tarar. 
    \item Quitar el tapón del matraz
    \item Descargar la bureta en el matraz aforado hasta que alcance la marca de 3ml
    \item Repetir 4 veces más
    \item Convertir cada masa a volumen corregido (mL) de acuerdo con la temperatura
    \item Realizar la diferencia de volúmenes
    \item Obtener un promedio de las masas
    \item Repetir este proceso con los siguientes volúmenes:
    \begin{itemize}
        \item 5
        \item 10
        \item 15
        \item 20
        \item 25
    \end{itemize}
    
    \item Calcular desviación estándar de los valores obtenidos
    \item Cotejar con la incertidumbre de la bureta
    \item Determinar si es apta para usarse o no
\end{enumerate}